\documentclass[english,twoside, openright, hidelinks, 11 pt, class=report,crop=false]{standalone}
\input{../../preamb}
\input{../bm}
\begin{document}

\grubr{opgalgsquare}
\begin{figure}
	\centering
	\subfloat[]{\includegraphics{\figp{opgalgsquare1_lf}}}\qquad \qquad
	\subfloat[]{\includegraphics{\figp{opgalgsquare2_lf}}}
\end{figure}
\abc{
\item Da det røde kvadratet har sidelengde $a$ og det grønne kvadratet har sidelengde $b$, har det største kvadratet sidelengde $a+b$. Hver av de hvite rektanglene har areal $ab$. Dermed har vi at
\begin{align*}
	A_{\text{størst kvadrat}}&=A_\text{rødt kvadrat}+A_\text{grønt kvadrat}+2A_\text{hvitt rektangel} \\
	(a+b)^2 &= a^2+b^2+2ab
\end{align*}
\item Det blå kvadratet har sidelengde $a-b$. Hver av de hvite rektanglene har areal $(a-b)b$. Dermed har vi at
\begin{align*}
	A_{\text{størst kvadrat}}&=A_\text{blått kvadrat}+A_\text{grønt kvadrat}+2A_\text{hvitt rektangel} \\
	a^2&= (a-b)^2+b^2+2(a-b)b \\
	a^2&= (a-b)^2+b^2+2ab-2b^2 \\
	a^2-2ab+b^2 &= (a-b)^2
\end{align*}
\item \begin{align*}
	a^2-b^2 &= A_{\text{størst kvadrat}}-A_\text{grønt kvadrat} \\
	&=A_\text{blått kvadrat}+2A_\text{hvitt rektangel} \\
	&=(a-b)^2+2(a-b)b \\ 
	&=(a-b)(a-b+2b) \\
	&=(a-b)(a+b)
\end{align*}
}
\newpage
\grubr{opgalgdeterminant}
\fig{determinant_lf}
Vi lar
\alg{
 E &= (a+c, 0) & F &= (0, b+d)
}
Med $ OE $ som grunnlinje har $ \triangle OEB $ høyde $ b $, altså er
\[  2A_{\triangle OEB}=(a+c)b \]
Tilsvarende er
\[  2A_{\triangle FDO}=(b+d)c \]
Da $ A_{\triangle OEB}=A_{\triangle CDF}$ og $ A_{\triangle FDO}=A_{\triangle EBC} $, har vi at
\alg{
	A_{\square ABCD}&=A_{\square OECF}-2A_{\triangle OBE}-2A_{\triangle FDO} \\
	&=(a+c)(b+d)-(a+c)b-(b+d)c \\
	&=(a+c)d-(b+d)c\\
	&= ad-bc
}

\grubr{opgalglikbrok} \\
Vi har at
\[ a= \frac{cb}{d} \]
Dermed er
\alg{
\frac{a-c}{b-d}=\frac{\frac{cb}{d}-c}{b-d}=\frac{c(b-d)}{d(b-d)}=\frac{c}{d}=\frac{a}{b}
}

\grubr{opgalgtverrsum}\\
Gitt et tall $n=abc $, hvor $ a $, $ b $ og $ c $ er sifrene til tallet 
Da har vi at
\alg{
n &= 100a+10b+c \\
&= 99a+99b+a+b+c
}
Leddene med 99 som faktor er delelige med 3, og dermed er $ n $ delelig med 3 hvis $ a+b+c $ er delelig med 3.
\end{document}


